\documentclass[aspectratio=169, 10pt, t]{beamer}

% === Metropolis Theme ===
\usetheme[
  progressbar=frametitle,
  numbering=fraction,
  block=fill,
  sectionpage=none
]{metropolis}

% === TUM Colors ===
\definecolor{tumblue}{HTML}{0065BD}
\definecolor{tumdark}{HTML}{003359}
\definecolor{tumsec}{HTML}{005293}
\definecolor{tumaccent}{HTML}{E37222}
\definecolor{tumgray}{HTML}{8C9AAF}
\definecolor{tumbg}{HTML}{F0F4F8}
\definecolor{tumlight}{HTML}{E8ECF0}
\definecolor{tumwarm}{HTML}{FEF6EE}
\definecolor{tumfaint}{HTML}{F8FAFB}

% === Apply TUM colors to Metropolis ===
\setbeamercolor{normal text}{fg=tumdark, bg=white}
\setbeamercolor{alerted text}{fg=tumaccent}
\setbeamercolor{example text}{fg=tumblue}
\setbeamercolor{frametitle}{fg=white, bg=tumblue}
\setbeamercolor{title separator}{fg=tumaccent}
\setbeamercolor{progress bar}{fg=tumaccent, bg=tumlight}
\setbeamercolor{progress bar in head/foot}{fg=tumaccent, bg=tumlight}
\setbeamercolor{block title}{fg=white, bg=tumblue}
\setbeamercolor{block body}{bg=tumbg}
\setbeamercolor{footline}{fg=tumgray}
\setbeamercolor{page number in head/foot}{fg=tumgray}
\setbeamercolor{item}{fg=tumblue}

% === Fonts ===
\setbeamerfont{title}{size=\Large, series=\bfseries}
\setbeamerfont{frametitle}{size=\large}

% === White math shorthand for frame titles ===
\newcommand{\tw}[1]{\textcolor{white}{$#1$}}

% === Footline ===
\setbeamertemplate{frame footer}{\color{tumgray}\scriptsize Ruben Kaiser\;\textbar\;Master Thesis Progress\;\textbar\;April 2026}

% === Packages ===
\usepackage[utf8]{inputenc}
\usepackage[T1]{fontenc}
\usepackage{amsmath, amssymb}
\usepackage{booktabs}
\usepackage{tabularx}
\usepackage{ragged2e}
\usepackage{xcolor}

% === Colored math commands ===
\newcommand{\gb}{\textcolor{tumblue}{g_{jk}}}
\newcommand{\gfn}{\textcolor{tumblue}{g}}
\newcommand{\Mb}{\textcolor{tumsec}{M}}
\newcommand{\Ub}{\textcolor{tumsec}{U}}
\newcommand{\Rb}[1]{\textcolor{tumsec}{R_{#1}}}
\newcommand{\Bb}{\textcolor{tumaccent}{B_{jk}}}

% === Subtle formula highlight (just a faint background strip) ===
\newcommand{\mathstrip}[1]{%
  \par\vspace{4pt}%
  \colorbox{tumbg}{\hspace{0.02\textwidth}\parbox{0.92\textwidth}{\vspace{6pt}%
    \centering #1%
  \vspace{6pt}}}\hspace{0.02\textwidth}%
  \par\vspace{4pt}%
}

% === Subtle side-accent for remarks (just a thin left rule) ===
\newcommand{\sidenote}[1]{%
  \par\vspace{4pt}%
  \hspace{2pt}{\color{tumgray}\vrule width 1.5pt}\hspace{8pt}%
  \parbox{\dimexpr\textwidth-14pt}{\footnotesize\color{tumgray}#1}%
  \par\vspace{2pt}%
}

% === Metadata ===
\title{Boundary Testing for Vision-Language Models}
\subtitle{Progress Update}
\author{Ruben Kaiser}
\institute{Supervisors: Oliver Karras, Andrea Stocco\;\textbar\;TUM}
\date{April 2026}

\begin{document}

% ============================================================
\maketitle
% ============================================================

% ============================================================
% SLIDE 1 — Decision System
% ============================================================
\begin{frame}{We describe systems that map inputs to categorical outputs.}

We consider any system that receives an input and produces one of finitely many outputs:

\mathstrip{$\displaystyle
  f \colon \textcolor{tumsec}{M} \to Y, \qquad |Y| \geq 2
$}

\vspace{2pt}

\begin{tabularx}{\textwidth}{@{}>{\bfseries}l Xl@{}}
\toprule
System & Input\; $m \in \textcolor{tumsec}{M}$ & \textbf{Output}\; $y \in Y$ \\
\midrule
Radiologist        & Chest X-ray        & Diagnosis \\
VLM (Qwen)         & Image + Prompt     & Classification \\
Credit scoring     & Application data   & Approve\,/\,Reject \\
\bottomrule
\end{tabularx}

\vspace{8pt}
\textit{\color{tumgray}All three share this structure. What follows applies to any such~$f$.}

\end{frame}

% ============================================================
% SLIDE 2 — Input Space
% ============================================================
\begin{frame}{The input space is structured and context-dependent.}

Each input dimension $j$ has a value space:\;
$S_j \in \bigl\{\mathbb{R},\; \mathbb{Z},\; \{c_1, \ldots, c_k\}\bigr\}$

\vspace{2pt}
The combinatorially complete space is\; $\Ub = \prod_{j=1}^{d} S_j$.

\vspace{4pt}
In a given context $C$ (dataset, prompt, deployment), only a subset is realized:

\mathstrip{$\displaystyle
  \textcolor{tumsec}{M} \;=\; \sigma(C) \;\subseteq\; \textcolor{tumsec}{U}
$}

\vspace{4pt}
\sidenote{A 2-class prompt and a 15-class prompt over the same images define different~$C$, therefore different~$\Mb$, therefore different boundary structures.}

\end{frame}

% ============================================================
% SLIDE 3 — Partition & Boundaries
% ============================================================
\begin{frame}{\tw{f} partitions \tw{M} into decision regions.}

$f$ assigns each $m \in \Mb$ to exactly one $y_k \in Y$:

\mathstrip{$\displaystyle
  \textcolor{tumsec}{R_k} = \bigl\{m \in \Mb \;\big|\; f(m) = y_k\bigr\}, \qquad \Mb = \bigsqcup_{k=1}^{K} \textcolor{tumsec}{R_k}
$}

\vspace{4pt}
Where two regions meet, a \textbf{boundary} exists.
For $|Y| \geq 2$ this is a structural consequence:
\[
  \partial \textcolor{tumsec}{R_{j}} \;\cap\; \partial \textcolor{tumsec}{R_{k}} \;\neq\; \emptyset \quad \text{for at least one pair}\; (j,\, k)
\]

\vspace{6pt}
If boundaries exist, the natural follow-up question is:

\vspace{2pt}
{\color{tumblue}\rule{0.05\textwidth}{1.5pt}}\\[4pt]
{\large\color{tumblue}\textit{What does their geometry look like?}}

\end{frame}

% ============================================================
% SLIDE 4 — Contrast Function g
% ============================================================
\begin{frame}{The contrast function \tw{g} makes boundaries measurable.}

For a class pair $(j, k)$, define:

\mathstrip{$\displaystyle
  \textcolor{tumblue}{g_{jk}}(m) \;=\; P(y_j \mid m) \;-\; P(y_k \mid m)
$}

\vspace{2pt}
\begin{columns}[T]
\begin{column}{0.42\textwidth}
\begin{align*}
  \textcolor{tumblue}{g} > 0  &\;\Rightarrow\; \text{prefers } y_j \\
  \textcolor{tumblue}{g} < 0  &\;\Rightarrow\; \text{prefers } y_k \\
  \textcolor{tumblue}{g} = 0  &\;\Rightarrow\; \text{\textbf{boundary}}
\end{align*}
\end{column}
\begin{column}{0.54\textwidth}
Properties:
\begin{itemize}\setlength\itemsep{2pt}
  \item $|\gb(m)|$\; = \;boundary proximity
  \item $\gb = -\textcolor{tumblue}{g_{kj}}$\; (antisymmetric)
\end{itemize}
\end{column}
\end{columns}

\vspace{3pt}
Boundary as zero-set:\quad $\textcolor{tumaccent}{B_{jk}} = \gfn^{-1}(0) = \{m \in \Mb \mid \gb(m) = 0\}$

\vspace{3pt}
\sidenote{$\gfn$ is the primary formal object.\; $\Bb$ is derivable from $\gfn$, but $\gfn$ carries strictly more information.}

\end{frame}

% ============================================================
% SLIDE 5 — Four Quantities
% ============================================================
\begin{frame}{Four geometric quantities follow from \tw{g} and \tw{\Delta g}.}

On the genotype lattice $G \subset \mathbb{Z}^d$, the lattice difference is \emph{exact}:\;
$\Delta_i\, \textcolor{tumblue}{g}(x) = \textcolor{tumblue}{g}(x + e_i) - \textcolor{tumblue}{g}(x)$

\vspace{6pt}

\begin{tabularx}{\textwidth}{@{}l >{\raggedright}p{3.6cm} X@{}}
\toprule
\textbf{Quantity} & \textbf{Definition} & \textbf{Measures} \\
\midrule
\textcolor{tumblue}{Margin}      & $|\gb(m)|$                                              & Distance to boundary in output space \\[3pt]
\textcolor{tumblue}{Sensitivity} & $\|\Delta \gfn(m)\| = \max_i |\Delta_i \gfn|$           & Responsivity to perturbation at $m$ \\[3pt]
\textcolor{tumblue}{Thickness}   & $|\gfn|\, / \,\|\Delta \gfn\|$                          & Perturbation budget until boundary \\[3pt]
\textcolor{tumblue}{Direction}   & $\arg\max_i |\Delta_i \gfn(m)|$                          & Most decision-relevant input dimension \\
\bottomrule
\end{tabularx}

\end{frame}

% ============================================================
% SLIDE 6 — Questions ↔ Quantities
% ============================================================
\begin{frame}{Each quantity answers a different question about the system.}

\begin{tabularx}{\textwidth}{@{}Xll@{}}
\toprule
\textbf{Question} & \textbf{Quantity} & \textbf{Formal} \\
\midrule
``How clear was this decision?''         & \textcolor{tumblue}{Margin}      & $|\gb(m)|$ \\[5pt]
``How fragile is this prediction?''      & \textcolor{tumblue}{Thickness}   & $|\gfn|\, / \,\|\Delta \gfn\|$ \\[5pt]
``Which input dimension matters most?''  & \textcolor{tumblue}{Direction}   & $\arg\max_i |\Delta_i \gfn|$ \\[5pt]
``How responsive is $\gfn$ here?''       & \textcolor{tumblue}{Sensitivity} & $\|\Delta \gfn(m)\|$ \\
\bottomrule
\end{tabularx}

\vspace{3pt}
All four derived from a single object:\; $\gb$.

\vspace{6pt}
\sidenote{These questions map independently to established concepts in systems psychology --- competence need {\scriptsize(Deci \& Ryan, 1985)}, perceived control {\scriptsize(Bandura, 1989)}, procedural justice {\scriptsize(Thibaut \& Walker, 1975)}.}

\end{frame}

% ============================================================
% SLIDE 7 — Literature Alignment (reframed)
% ============================================================
\begin{frame}{Existing work addresses the same underlying structure.}

\vspace{2pt}
{\small\itshape\color{tumdark}%
``Deep learning systems are increasingly deployed in safety- and security-critical domains [\,\ldots\,] where the correctness and predictability of a system's behavior for corner case inputs are of great importance.''}

\hfill\textcolor{tumgray}{\footnotesize--- Pei et al.\ (DeepXplore), SOSP 2017}

\vspace{10pt}
These concepts map directly onto $\gb$:

\vspace{4pt}
\begin{tabularx}{\textwidth}{@{}lX@{}}
``corner case inputs''   & $\longleftrightarrow\;$ inputs near $\Bb$\; {\small(low \textcolor{tumblue}{margin})} \\[5pt]
``correctness''          & $\longleftrightarrow\;$ whether $f(m)$ is stable\; {\small(\textcolor{tumblue}{thickness})} \\[5pt]
``predictability''       & $\longleftrightarrow\;$ knowing $\gfn(m)$ before observing $f(m)$ \\
\end{tabularx}

\vspace{10pt}
$\gb$ provides a shared formal language for these intuitions.

\end{frame}

% ============================================================
% SLIDE 8 — BT Definition
% ============================================================
\begin{frame}{For opaque systems, \tw{g} is only empirically accessible.}

\vspace{4pt}
\begin{columns}[T]
\begin{column}{0.50\textwidth}
\begin{align*}
  \textbf{Transparent}\; f  &\;\Rightarrow\; \gfn \text{ via inspection} \\[6pt]
  \textbf{Opaque}\; f       &\;\Rightarrow\; \gfn \text{ via queries only}
\end{align*}
\small (DNNs, VLMs, any black-box system)
\end{column}
\begin{column}{0.46\textwidth}
\mathstrip{$\displaystyle
  \textcolor{tumaccent}{BT}(f) \;:=\;
  \text{\small systematic empirical approximation of\;} \gb
$}
\end{column}
\end{columns}

\vspace{4pt}

\begin{tabularx}{\textwidth}{@{}lX@{}}
\toprule
\textbf{Approach} & \textbf{Characterization} \\
\midrule
XAI (LIME, SHAP)            & Local, post-hoc --- approximates $\gfn$ at single points \\[3pt]
\textcolor{tumaccent}{Boundary Testing}  & Global, systematic --- maps the $\gfn$-landscape \\
\bottomrule
\end{tabularx}

\vspace{6pt}
\sidenote{VLMs are the hardest case --- but not the reason for the framework.}

\end{frame}

% ============================================================
% SLIDE A — From Geometry to Test Decisions
% ============================================================
\begin{frame}{From boundary geometry to test decisions.}

Each $\gfn$-quantity implies a concrete testing parameter:

\vspace{2pt}

\begin{tabularx}{\textwidth}{@{}l >{\raggedright}p{3.2cm} X@{}}
\toprule
\textbf{Quantity} & \textbf{Test Decision} & \textbf{Rationale} \\
\midrule
\textcolor{tumblue}{Thickness}   & Priority\; {\small(which pairs first?)}   & Fragile boundaries need more coverage \\[2pt]
\textcolor{tumblue}{Direction}   & Modality\; {\small(text or image?)}       & Test the dimension that controls the boundary \\[2pt]
\textcolor{tumblue}{Sensitivity} & Budget\; {\small(how many seeds?)}         & High-sensitivity regions change fast \\[2pt]
\textcolor{tumblue}{Margin}      & Stopping\; {\small(when conclusive?)}      & Low margin $=$ near boundary $=$ informative \\
\bottomrule
\end{tabularx}

\vspace{4pt}
\begin{center}
\textcolor{tumblue}{Boundary geometry} $\;\xrightarrow{\;\; \gb \;\;}\;$ \textcolor{tumblue}{Test parameters}
\end{center}

\end{frame}

% ============================================================
% SLIDE B — Cross-Model & Cross-Class Hypotheses
% ============================================================
\begin{frame}{Two hypotheses that the framework makes testable.}

\textbf{H1: Boundary geometry is a property of the task, not the model.}

\vspace{2pt}
\sidenote{Run the same seeds on a second model (e.g.\ LLaVA alongside Qwen). If \textcolor{tumblue}{thickness} rankings correlate across models $\Rightarrow$ the task determines fragility. If they diverge $\Rightarrow$ the architecture does.}

\vspace{10pt}

\textbf{H2: Boundary fingerprints cluster by semantic similarity.}

\vspace{2pt}
\sidenote{The activation patterns per class pair form a boundary fingerprint. Do these cluster taxonomically --- birds vs.\ reptiles vs.\ fish? If yes, boundary structure reflects semantic structure, not just statistical proximity.}

\vspace{14pt}

Both are testable with the existing pipeline + one additional model run.

\vspace{4pt}
\sidenote{A \textbf{text--image dominance ratio} per class pair --- which modality controls the boundary --- would be a VLM-specific contribution beyond classical BVA.}

\end{frame}

% ============================================================
% SLIDE C — Running & Planned Experiments
% ============================================================
\begin{frame}{Running and planned experiments.}

\vspace{2pt}

\begin{tabularx}{\textwidth}{@{}>{\raggedright}p{3.8cm} l X@{}}
\toprule
\textbf{Experiment} & \textbf{Status} & \textbf{Goal} \\
\midrule
SMOO 2-class (55 seeds)          & \textcolor{tumblue}{done}           & Baseline boundary geometry \\[3pt]
PDQ 15-class (6 seeds)           & \textcolor{tumblue}{done}           & Multi-class boundaries, two stages \\[3pt]
v2\_strategies                    & \textcolor{tumaccent}{running}      & Extended Stage-1 strategy coverage \\[3pt]
v2\_deep                          & \textcolor{tumaccent}{running}      & Deeper optimization per seed \\[3pt]
v2\_gap                           & \textcolor{tumaccent}{running}      & Targeted low-margin class pairs \\[3pt]
Cross-model (LLaVA)              & {\color{tumgray}planned}           & H1: task vs.\ model property \\[3pt]
Taxonomic clustering             & {\color{tumgray}planned}           & H2: semantic boundary structure \\
\bottomrule
\end{tabularx}

\vspace{10pt}
\sidenote{$\gfn$ as a monitoring signal --- periodically sampling boundary geometry across model versions --- is a natural deployment extension, sketched in the discussion.}

\end{frame}

\end{document}
